\documentclass[11pt,a4paper]{article}
\usepackage[top=42pt,bottom=42pt,left=70pt,right=70pt]{geometry}  % to make pages wider

\usepackage{amsmath}
\usepackage{amsfonts}
\usepackage{graphicx}  % for figures etc.
\usepackage{url}  % to handle URLs 
\usepackage[comma,authoryear]{natbib}  % for author-date ref style --  allows \citet (textual), and \citep (parenthesized)
\usepackage[hidelinks]{hyperref}
\hypersetup{allcolors=[rgb]{0, 0, 0.33},colorlinks=true}

\renewcommand {\cite} {\citep}  % default for cite is citet in natbib - changes it to citep
\renewcommand\bibname{References}    % if not using newrucsthesis sty file

\usepackage{booktabs}
\usepackage{array}
\usepackage{enumitem}
\usepackage{amssymb}
\begin{document}

	
% Create a title
\title{Literatue review }
\author{ Karabo Ntsie }
\date {May 2026}
\maketitle

\section{Background Concepts}

\section{The African NLP landscape}
This section discusses the African NLP landscape and the marginalisation of African languages in NLP research. This section does this by examining the measurable scale of underrepresentation across resource tiers, the documented quality of wen-crawled multilingual data, the limitations of transfer learning and data augmentation as responses to data scarcity, and the underperformance of large multilingual models on African languages despite their scale. Together, these subsections establish the challenge faced by African NLP. It is within this broader context that the specific challenges of Sesotho corpus construction, examined in the following section, must be understood.

\subsection{The quantitative scale of under representation}
The underrepresentation of African languages in NLP is not only observable but measurable through resource audits and corpus analysis. \citet{joshi2020state} classify the world’s languages into six tiers based on the availability of labelled and unlabelled data. Their findings show that 2,191 languages, representing 88.38\% of all languages, fall into the lowest tier with virtually no digital resources. \citet{adebara2022towards} extend this analysis to African languages and report that 542 fall into this lowest category, 26 have minimal resources, and only one language, Afrikaans, reaches a higher tier due to its web presence. No Indigenous African language appears in the top tiers, highlighting the scale of the imbalance.

However, counting datasets alone does not capture the full extent of the problem. \citet{yu2022beyond} analyse 156 multilingual datasets covering 222 languages and find that 68\% of these languages have no manually annotated data. Among low-resource languages, 84.9\% of datasets rely on automatically generated labels, which are often noisy and unreliable. Only 40 languages, representing 18\% of those surveyed, include text written specifically by humans for NLP tasks. Most datasets rely on limited sources such as Wikipedia, which does not reflect the full range of language use. The study further shows that dataset availability strongly correlates with the number of NLP researchers who speak the language ($\rho$ = 0.71) and with the availability of crowdworkers ($\rho$ = 0.58–0.59), indicating that participation, rather than purely technical constraints, drives the resource gap.

These disparities have direct consequences for model performance. The gap between high- and low-resource languages spans several orders of magnitude. English has over 6 million Wikipedia articles and hundreds of labelled datasets, while low-tier languages may have only a few pages and no standardised resources. \citet{hedderich2021survey} show that English supports at least eighteen NLP task categories, while languages such as Nahuatl and Quechua support fewer than eight, with entire capability levels absent.

Efforts to address this gap demonstrate both progress and limitations. \citet{nekoto2020participatory} shows that community-driven work can produce meaningful results through the Masakhane project, which created 45 machine translation benchmarks across 32 African languages, including 16 classified as left behind and 11 as scraping by. For many of these languages, this was the first instance of human evaluation. However, the process required sustained coordination of over 400 participants across at least 20 countries, highlighting the difficulty of scaling such efforts and the need for more automated approaches.

The architecture of modern NLP further reinforces this imbalance. Neural models require large volumes of unlabelled data, which low-resource languages lack. \citet{joshi2020state} argue that this creates a cycle in which methods designed to reduce dependence on labelled data still depend on resources that many languages do not have. \citet{yu2022beyond} show that cross-lingual datasets attempt to address this through translation, yet only 33 of the 156 datasets surveyed follow this approach, and 61.9\% of these rely on automatic translation. This introduces artefacts that do not reflect native linguistic use, reducing data quality.

Data quality issues further undermine existing resources. \citet{adebara2022towards} report that among 205 multilingual datasets containing African languages, at least 15 were completely erroneous, many contained less than 50\% correct data, and 82 were mislabelled or used ambiguous language codes. In the Flores-101 dataset for Yorùbá, 12.4\% of entries contain incorrect tone marks, 5.29\% contain spelling errors, and 2.7\% show inconsistent spelling. In tone languages, such errors can change meaning and negatively affect model training. Yu et al. (2022) confirm that 53 out of 156 datasets rely on automatically induced labels and often lack validation measures, such as inter-annotator agreement.

This imbalance is also reflected in research output. \citet{joshi2020state} show that high-resource languages dominate major NLP conferences, while low-resource languages receive significantly less attention. \citet{yu2022beyond} further shows that although 62 datasets were created for multilingual tasks covering 217 languages, only 16 languages have dedicated monolingual benchmarks, and these are primarily languages with established research communities. Multilingual datasets are more likely to rely on translated or automatically generated data, which is typically lower in quality.

Taken together, these findings show that the challenges facing African NLP form a structural cycle. Limited high-quality data constraints model development, reducing incentives to generate new data. \citet{nekoto2020participatory} link this to broader societal factors, while \citet{yu2022beyond} describes it as a people gap, in which resource availability depends on the presence of language-proficient researchers. Addressing this issue requires scalable and systematic approaches to data collection and curation, alongside efforts to support African NLP communities.

\subsection{The quality failure of web-crawled multilingual data}

The prevailing assumption that web-crawled corpora, supplemented by curated reference sources such as Wikipedia, provide an adequate foundation for multilingual NLP rests on a circular dependency that the literature has only recently begun to interrogate systematically. Wikipedia is routinely used as a high-quality reference to train the language identification classifiers that partition Common Crawl along language boundaries. However, the quality of Wikipedia itself, particularly for non-English languages, is far from assured. \citet{tatariya2024good} empirically demonstrates that MinHash deduplication alone affects over 28\% of all non-English Wikipedia articles. In contrast, \citet{perelkiewicz2024review} reports that approximately 40\% of the Common Crawl archives consist of exact duplicates, with paragraph-level duplication reaching 70\% in some snapshots. The field is therefore not operating on a clean spectrum from noisy web data to curated reference material, but within an ecosystem in which quality failures are mutually reinforcing and systemically reproduced across the entire data pipeline.

These failures are not evenly distributed between languages. \citet{tatariya2024good} four-tier quality taxonomy of non-English Wikipedias reveals that the languages most severely affected are disproportionately from the Global South: Yorùbá Wikipedia, for instance, loses nearly 60\% of its articles to exact-match deduplication, with many removed entries consisting of single placeholder tokens rather than substantive linguistic content. The errors introduced at this stage propagate downstream, since the compromised Wikipedia data used to train language identification classifiers then distorts the composition of the corpora that those classifiers are meant to organise. \citet{kwon2026afroscope} identify domain skew as a further compounding factor, noting that available African language data is heavily concentrated in religious texts and Bible translations, failing to capture the register diversity necessary for robust language modelling. This is reinforced by the geolocation bias documented by \citet{perelkiewicz2024review}, whereby countries with large African-language-speaking populations remain radically underrepresented in the Common Crawl relative to their demographic weight, meaning that web-crawled data attributed to these languages frequently originates from narrow and unrepresentative community segments.

A further dimension of quality failure operates at the level of linguistic structure itself, largely invisible to the document-level filters dominant in the field. \citet{de2007automatic} demonstrates that the digitisation of African language corpora through OCR processing and non-standard encodings systematically strips diacritical marking, silently removing phonological and morphological information encoded in characters outside the standard Latin alphabet. Crucially, this degradation evades the heuristics most commonly applied in large-scale corpus construction — length thresholds, type-token ratios, entropy metrics — all calibrated against assumptions derived from high-resource, English-language contexts. A document stripped of its diacritics will pass every standard quality filter yet remain linguistically impoverished in ways that undermine any downstream task dependent on phonemic or morphological distinctions. This points to a fundamental mismatch between the quality frameworks the field has developed and the specific forms of degradation most prevalent in African language data, a mismatch that cannot be resolved by applying existing pipelines more aggressively or at greater scale. \citet{tatariya2024good} recommendation for per-language quality assessment, \citet{kwon2026afroscope} emphasis on domain diversity, and \citet{de2007automatic} demonstration of sub-lexical degradation collectively converge on the same diagnosis: the quality failure of web-crawled multilingual data is structurally reproduced within the architecture of the data pipeline itself, and its consequences fall with disproportionate severity on African and other low-resource languages.

\subsection{The limitations of transfer learning and augmentation as solutions}
\citet{hedderich2021survey} catalogue an extensive array of such methods, yet note that most assume auxiliary resources — parallel corpora, bilingual dictionaries, or machine translation systems — that do not exist for truly low-resource languages. This section argues that these technical fixes circumvent rather than solve the problem of data quality, revealing a gap that automated curation is positioned to fill. 

Cross-lingual transfer learning, whether data- or model-based, relies on prerequisites that are rarely met in African languages. \citet{hedderich2021survey} identify three dimensions of resource availability — task-specific labels, unlabeled text, and auxiliary data — and demonstrate that most transfer techniques require at least two. For languages such as Tiv, \citet{achir2026preservation} document a "doubly low-resource scenario": no parallel corpus, no POS tagset, no MT system, and no ASR resource. Where resources exist, performance remains poor. \citet{garcia2025cross} report that zero-shot transfer on MasakhaNER 2.0 yields only 24.9\% F1 for isiXhosa and 43.9\% for isiZulu — a fraction of the 92.4\% achieved for English. \citet{thangaraj2024cross} further shows that fine-tuning multilingual models on Kirundi results in 74\% catastrophic forgetting on the original Kinyarwanda, trading off source performance for modest target gains.

Data augmentation techniques — back-translation, switchout, and sentence concatenation — operate on existing data rather than generating higher-quality data. \citet{lamar2023measuring} introduces the Data Augmentation Advantage (DAA) metric, quantifying how many true pairs a synthetic pair replaces. For English-Italian at 500K training pairs, sentence concatenation achieved DAA of 1.38 (each synthetic pair worth 1.38 true pairs). However, at 1M pairs, the same method yielded -0.03, which is actively harmful. For true low-resource languages (English-Romany, 24K pairs), DAA varied unpredictably from -0.13 to 3.0 across augmentation types, demonstrating that benefits are neither consistent nor guaranteed.

\citet{oduwole2025scarcity} extend this analysis to six African language pairs. While augmentation improved BLEU by up to 25\% relative to the baseline, absolute gains remained marginal. English-Hausa improved from 5.10 to 9.65 BLEU; English-Yoruba from 6.42 to 9.14 BLEU. For French-Fon, the best augmentation achieved only 3.45 BLEU (baseline 0.89) — functionally useless for real-world translation. Critically, English-Swahili — the best-resourced pair with 30,782 parallel sentences — showed no improvement, with the baseline (25.80 BLEU) outperforming all augmented variants. Augmentation cannot compensate for inadequate seed data; it may even harm when sufficient data already exists.

For ASR, \citet{bartelds2023making} reports that self-training on four minority languages yields relative WER reductions of 6.3–13.9\% when starting from just 24 minutes of transcribed speech. However, diminishing returns set in rapidly: adding 672 minutes of self-trained data produced only a 20.5\% relative reduction — less than double the gain from the first 168 minutes. The quality of generated transcriptions is constrained by the teacher model, which itself is trained on limited seed data.

LLM-based synthetic data generation promises scalability but introduces new failure modes. \citet{gyamfi2026synthetic}, generating 50,000 Swahili sentiment samples with LLM judging, achieved 80.95\% agreement with human ground truth. However, they identify "translationese" — grammatically correct but culturally hollow text — as a persistent problem. Generated proverbs, while linguistically accurate, were frequently incongruent with context, revealing that LLMs approximate statistical patterns without pragmatic grounding.

\citet{Zhou2024} report more successful automation for medical records: their OCR-NLP pipeline achieved 98.66\% accuracy (vs 98.65\% manual) while reducing collection time from 63.6 to 3.6 minutes per patient. However, the authors caution that this success depends on highly constrained documents with "white background with black text" — conditions that do not hold for heterogeneous web data.\citet{oduwole2025scarcity} similarly observe that augmentation techniques degrade as languages become more resource-scarce: for French-Wolof, even the best augmentation (7.70 BLEU) remains far below usable quality. Augmented data inherits the biases and limitations of its source.

The limitations surveyed above point to a deeper structural issue. \citet{hedderich2021survey} acknowledge that most low-resource techniques require auxiliary resources — MT systems, bilingual dictionaries, parallel corpora — that are absent for languages like Tiv. \citet{yu2022beyond} demonstrate that the availability of manually annotated datasets correlates with the number of language-proficient researchers ($\rho$ = 0.71). The scarcity is primarily human, not technical. Cross-lingual transfer and data augmentation are responses to this scarcity, but they function as palliatives. \citet{garcia2025cross} concluded that there is no universal, multilingual pretrained model for every language and domain. These limitations justify a complementary approach: scalable, quality-first data curation that addresses the source of scarcity rather than its symptoms.

\subsection{The failure of large multilingual models for African languages}

Despite being positioned as solutions for low-resource languages, large multilingual models such as mBERT and XLM-R underperform on African languages. \citet{hussen2025state} quantifies the gap: African languages represent 28.57\% of the world's 7000 languages, while multilingual models allocate fewer than 10\% of language slots. mBERT supports 104 languages, but only 6 are African; a proportional representation would require approximately 30. mT5 supports 101 languages, including 14 African languages, whereas it should support 29. XLM-R supports 100 languages, with 8 African languages. Across all surveyed models, 42 of over 2000 African languages receive support. Amharic, Swahili, Afrikaans and Malagasy dominate the coverage, while 98\% of African languages have no model support.

\citet{chang2024multilinguality} explains the curse of multilinguality, where adding more languages to a fixed-capacity multilingual model systematically degrades performance per language. For high-resource languages, adding 1 billion multilingual tokens is equivalent to removing 63-88\% of the target language's monolingual data. For low-resource languages, the benefits from multilingual data plateau quickly, leading to too much data and eventually hurting performance as models reach their limits. Performance gains depend on the syntactic similarity of added languages. This curse explain the under representation documented by \citet{hussen2025state} where 42 of over 2000 African languages face additional barrier that \citet{chang2024multilinguality}'s analysis identifies as critical to performance: 3 of 23 active scripts are supported =, making most African languages unseen scripts that partially script-agnostic models cannot handle; tokenisaton is biased toward Latin-script languages; and digitised corpus availability is severely limited. Where \citet{chang2024multilinguality} shows that monolingual tokenisers achieve 93.7\% vocabulary coverage with minimal data, \citet{hussen2025state} shows that multilingual tokenisers fail on African scripts entirely, thus producing worse outcomes for African languages than for low-resource languages with compatible scripts and digitised corpora.

\citet{ogueji2022afriberta} directly tests the generic multilingual approach against a targeted alternative. AfriBERTa, trained on 0.94GB of curated African language data with 108 million tokens, matches or exceeds XLM-R and mBERT on Hausa, Amharic and Swahili despite being 3 times smaller and trained on 100 times less data. It performs well on languages it was never pretrained on (Luo, Wolof, and Luganda). In text classification, AfriBERTa outperforms XLM-R and mBERT by over 10 F1 points on Yoruba and 7 points on Hausa, thus concluding that pretraining on similar low-resource languages may be better than mixing high- and low-resource languages, which contradicts the assumption underlying generic multilingual models.

\citet{rahman2026pashtocorp} replicates this finding for Pashto. XLM-R bases achieve only 19\% F1 on WikiANN Pashto NER. Continued pretraining on PashtoCorp, a targeted 1.25-billion-word corpus, improves F1 by 10\% relative to 21\% and reduces training variance nearly 7 times from $\pm 4.7\%$ to $\pm 0.7\%$. The gain is largest at 50 training sentences ($\pm 27\%$), demonstrating that targeted data most benefits the lowest-resource scenarios.

\citet{brown2025pula} shows the ceiling of targeted approaches. Pula, a bilingual Sestswana-English model trained on curated data, outperforms GPT-4o and Gemini 1.5 Pro on English-Setswana translation tasks and exceeds Llama 3.1 70B on all Setswana reasoning benchmarks despite having far fewer parameters.

Generic multilingual models fail African languages because of their design assumptions with maximising language count, using web-crawled noisy data, and script-agnostic tokenisation, which conflict with the characteristics of African languages. African-centric models trained on curated, script-appropriate data with linguistically similar language combinations consistently outperform larger generic models. The problem is an incorrect design, not insufficient scale.


\section{The Sesotho context}
This section discusses the Sesotho language and its standing in the natural language processing landscape. Understanding this standing requires first situating Sesotho within its linguistic, demographic, and sociopolitical context, and then examining how these factors intersect with the resource requirements of contemporary NLP.

The structure of this section proceeds from Sesotho's linguistic profile to the orthographic split, and ends with Sesotho's digital presence and the critical gaps in current corpora.

\subsection{The linguistic profile of Sesotho}
Sesotho is a Bantu language of the Sotho-Tswana group, spoken predominantly in Lesotho and South Africa, with over 13 million speakers, of whom 5.6 million are first-language speakers and 7.9 million have Sesotho as a second language \citep{mokoaleli2020sesotho}. The scale of speakers establishes Sesotho as a linguistically significant community that deserves computational support. Its continued marginalisation in human language technologies is therefore a misalignment between its linguistic structure and the underlying assumptions of current NLP methods.

This misalignment is rooted in its typological profile. Within the broader framework of morphological typology, Sesotho belongs to the agglutinative branch of the synthetic languages, a classification it shares with the Bantu group, and it stands in direct contrast to analytic or isolating languages such as English \citep{kambarami2021out,masethe2022word}. In agglutinative languages, words are constructed with morphemes, each carrying a single segmentable meaning \citep{kambarami2021out}. However, in Sesotho, this morphological organisation challenges token-based and vocabulary-driven computational methods, as meaning is not distributed across independent units but is encoded through tightly interdependent morphological structures that cannot be easily segmented. 

The noun class system exemplifies these challenges. Each noun class is distinguished by unique prefixes which govern concordial agreement across the entire sentence, determining the form of subject markers, qualificative particles and associative markers on all co-occuring constituents \citep{mokoaleli2020sesotho}. This creates dense, non-local dependencies that are poorly captured by models assuming token independence. The situation is further complicated by prefix drop, which introduces variability in surface forms and obscures underlying grammatical structure, undermining approaches that rely on stable orthographic cues.

Sesotho's morphology is generative because of a productive derivational process, alongside emerging patterns such as the extension of the prefix /bo-/ across multiple lexical categories, which demonstrate that the language continually produces novel, well-formed expressions \citep{mokoaleli2020sesotho,moloi2014morphology}. This open-ended productivity renders finite lexicons incomplete and exposes the limitation of closed-vocabulary NLP systems: their inability to generalise beyond observed forms.

Comparative evidence from related Sotho-Tswana languages reinforces that these challenges are systemic. Languages such as Sesotho sa Leboa (Sepedi) exhibit similar patterns of layered affixation and context-dependent polysemy, reflecting a shared agglutinative structure that generates dense and flexible form-mappings \citep{masethe2022word}. The computational consequences of agglutinative morphology lead to persistent out-of-vocabulary (OVO) failures, as productive affixation generates valid but unseen forms \citep{kambarami2021out}. While Sesotho's disjunctive orthography distributes morphological information across tokens, in contrast to conjunctive systems such as Shona, the underlying issue remains: models that treat words as stable, enumerable units fail to capture the combinatorial nature of morphology.

These observations highlight the limitations of prevailing NLP paradigms for Sesotho. Its morphological richness introduces an added complexity and directly contradicts the assumptions underpinning tokenisation, vocabulary design, and generalisation. Therefore, effective computational approaches must operate at the sub-word level and model morphological processes, or Sesotho's linguistic richness will manifest as a systemic failure in existing NLP systems.

\subsection{The orthographic split}
Sesotho's linguistic richness does not end at its morphological richness but extends to its challenge of having two competing orthographies. French missionaries developed the Lesotho orthography in the 19th century, while the South African orthography was formalised in 1959 \citep{makutoane2022people}. Harmonisation attempts began as early as 1927 and failed repeatedly because Lesotho authorities objected to the proposed changes \citep{matlosa2017sesotho}. The orthographic split was sustained by the geographical distribution of the Basotho nation spanning two countries; therefore, Sesotho text is being produced simultaneously by two populations under two different orthographies. Therefore, web scrapers pulling data for Sesotho text inevitably pull documents from both sides because South African Sesotho (SAS) and Lesothoan Sesotho (LS) superficially look similar, leading naive language identification to treat them as one variety. Thus, the competing orthographies will remain a permanent feature of the Sesotho digital landscape and must be accommodated. 

The SAS/LS split operates at two levels of the writing system simultaneously, each with distinct implications for automated processing. At the graphemic level, the divergence is systematic and rule-governed: SAS substitutes kg, tjh, tsh and d where LS writes kh, ch, tš and l respectively, omits the apostrophes LS uses to mark contracted sounds, drops the vowel diacritics LS preserves, and renders semi-vowels as ya and wa where LS writes ea and oa \citep{makutoane2022people}. The consistency of these substitutions is computationally significant because the pattern is rule-governed rather than arbitrary, and the variants are distinguishable at the character level, providing reliable anchor features for an orthographic classifier without requiring semantic analysis. \citet{motaung2025sesoda} confirm this characterisation independently in a computational dataset construction context, documenting the same consonant cluster divergence as a systematic feature requiring explicit handling rather than normalisation. Their morphological profile of Sesotho, with an out-of-vocabulary rate of 73\%, 18 noun classes, and an agglutinative index of 4.7, is substantially higher than Swahili (41\% OOV) and Zulu (58\% OOV), establishing that these graphemic differences interact with an exceptionally productive morphological system, meaning a misclassified orthographic variant does not produce a localised error but propagates inconsistency across morpheme boundaries throughout the document. The diacritic omission in SAS deepens this problem at the phonological level. \citet{matlosa2017sesotho} identifies the consequence as systematic underspecification: in SAS, the single grapheme "e" encodes three distinct phonological values and "o" encodes three others. For any downstream NLP task depending on lexical disambiguation or morphological analysis, two documents, one from each orthographic tradition, may therefore be statistically incomparable even when describing the same content, because the character distributions a model learns from are structurally different. A pipeline that conflates the two variants does not merely introduce noise; it produces a corpus whose internal statistical properties are incoherent.

\citet{lin-etal-2024-modeling} demonstrated that modes trained on orthographically diverse data with explicit variant labelling outperform those trained on normalised data in Nigerian Pidgin, a language without a unified writing convention. Therefore, the orthographic split is structured linguistic information that can be leveraged to improve downstream NLP performance.

\subsection{Sesotho's digital presence and corpus gaps}
The development of NLP tools for Sesotho is fundamentally constrained by the scarcity of high-quality supervised data for advanced computational tasks \citep{sibeko2022overview}. While Sesotho has over 13 million speakers across South Africa, Lesotho and Zimbabwe, it remains classified as a low-resource language \citep{mokhosi2024sesotho}. This classification reflects the language's digital representation. The literature reveals the abundance of certain categories of text due to specific institutional or religious initiatives. However, the practical utility is frequently undermined by orthographic inconsistency and a dearth of linguistically diverse, original material.

Of the literature surveyed, the digitally available Sesotho can be categorised into three distinct streams of translated government documents, religious/parallel texts, and emerging social datasets.

The most substantial and professionally curated source of monolingual data derives from South African government initiatives. The Autshumato Machine Translation project and the National Centre for Human Language Technology (NCHLT) text collection constitute the backbone of available raw text. \citet{Sibeko2024} reports utilising approximately 4.6 million from the Authshumato Sesotho dataset and a further million tokens from the NCHLT collection. These corpora are particularly valuable because they consist of texts professionally translated from English, thereby ensuring a relatively high standard of linguistic accuracy and adherence to standardised government terminology. \citet{marivate2020investigating} similarly leveraged these resources to build vectorisers for Sotho-Tswana languages, noting their utility in overcoming the "small data" problem inherent to low-resource languages (LRLs).

The second stream, being religious text specifically from the 1960 Sesotho Bible and Jehovah's Witness publications (JW300 corpus), represents a significant source of parallel and monolingual data. \citet{agic2019jw300} introduced JW300, a massive parallel corpus covering 300 languages, which includes substantial Sesotho alignments. Although there may be an ideological bias, acknowledged by \citet{agic2019jw300}, the corpus provided a rare source of sentence-aligned, cross-lingual information. \citet{Sibeko2024} utilised nearly one million tokens from the Sesotho Bible, arguing that biblical text remains a culturally embedded and orthographically stable resource.

The third stream is from emerging social media and online news. \citet{marivate2020investigating} harvested news headlines from the SABC Facebook page (Motsweding FM) in Setswana and Sepedi, establishing a methodology applicable to Sesotho. More concretely, \citet{mokhosi2024sesotho} manually scraped and curated the Sesotho News (SN) dataset from two online newspapers using Lesotho orthography, yielding 2401 headlines annotated for sentiment analysis. In the largest aggregation to date, \citet{ronny2025advancing} expanded the SAfriSenti corpus, which includes Sesotho as a part of a multilingual Twitter dataset, thus marking a critical shift toward capturing colloquialisms and code-switching.

Despite the availability of these three streams of data, the literature consistently identifies orthographic inconsistency as the primary quality barrier to aggregating them into larger, more powerful corpora. Systematic differences between the South African Sesotho (SAS) and Lesotho Sesotho (LS) orthographies affect fundamental text properties. Texts written to the two standards cannot be combined straightforwardly without extensive normalisation, a labour-intensive process that requires manual verification. Consequently, resources developed for one orthographic standard effectively exclude the linguistic community that uses the other orthographic standard. As noted by \citet{mokhosi2024sesotho}, Lesotho's orthography of Sesotho is largely sidelined from mainstream NLP research, limiting the availability of language-specific services for 1.85 million speakers. This orthographic fragmentation, combined with the domain-specific nature of each data stream, means that Sesotho NLP researchers often face a difficult choice: they must either work within the constraints of a single orthographic standard and domain, sacrificing broader representativeness, or invest prohibitive resources in manual normalisation to create a more comprehensive but smaller curated dataset. 

\section{Data collection and data quality}
The preceding sections established that African languages face a structural data crisis rooted not only in scarcity but in the poor quality of what exists, and that Sesotho presents a specific combination of orthographic and morphological challenges that generic solutions cannot address. This section examines the technical literature that informs how those challenges can be addressed in practice. It proceeds in three parts. First, it establishes what data quality means in the context of a low-resource corpus and how it can be operationalised into measurable criteria. Second, it surveys existing automated pipeline approaches and identifies where they succeed and where they fail for languages like Sesotho. Third, it examines language identification as a discrete pipeline component, with particular attention to the asymmetry between false positives and false negatives that makes LID design decisions consequential for corpus integrity. Together, these subsections provide the technical grounding for the pipeline design proposed in this project and establish that the gap to be filled lies not in any single component but in their principled integration.
\subsection{Defining and measuring data quality in NLP}
The quality of a text corpus is a measure of its fitness for use given the specific research context. For a low-resource language such as Sesotho , where digitised text is scarce and web-sourced data contains unknown noise, establishing a rigorous definition of quality is a prerequisite to downstream NLP tasks. \citet{hurtado2022documents} propose a total corpus quality framework that systematically decomposes error sources into three sources: source errors (inherent in original material, such as damaged documents or low-quality scans), and research inference errors (stemming from sampling, coverage, or specification decisions). For the Sesotho pipeline, source errors arise during OCR processing and sentence extraction, and research inference errors relate to sampling strategy for human evaluation and specification of what constitutes as Sesotho.

\citet{wang2026modernizing} extends this analysis by identifying four recurring challenges specific to unstructured data which are format heterogeneity (data arriving in PDFs,HTML and scanned images), semantic ambiguity (the same word having different meanings in different contexts), syntactic variability (spelling variations, code-switching, and informal register), scale (the impossibility of manual inspection for large volumes). Each of these challenges are relevant to the Sesotho language as format heterogeneity arises fro different formats from diverse web sources, syntactic variability due to the coexistence of South African and Lesotho orthographic standards, and scale demands  automated quality assessment. \citet{xiao2023evaluating}  introduce METRICEVAL, a framework that formalises two concepts for evaluating metrics which are reliability (consistent results across repeated measurements) and validity (measuring what is intended). Although METRICEVAL was designed to evaluate NLG (natural language generation) metrics rather than raw corpora, its distinction between reliability and validity provides a useful insight that a corpus has high reliability if repeated sampling yeilds consistent quality estimates, and high construct validity if quality criteria genuinely reflect usable Sesotho text.

Combining these frameworks, this mini-thesis operationalises corpus quality through three measurable criteria. First, language purity which addresses  the basic level of fitness of use which is how much of the corpus is actual Sesotho. Following \citet{kreutzer-etal-2022-quality} a random sample of 100 sentences will be annotated using binary taxonomy with a target of 90\% correct sentences. This draws on \citet{hurtado2022documents} concept of textual representation error.  Second, orthographic awareness measures the pipeline's ability to correctly classify South African versus Lesotho Sesotho variants on a held-out test set of 200 labelled sentences, with a target F1-score of at least 0.85. This criterion directly operationalises \citet{wang2026modernizing}'s syntactic variability challenge. Third, downstream utility demonstrates that quality improvements translate into practical benefit: a model trained on pipeline-curated data must achieve a statistically significantly higher F1-score (p < 0.05) on a text classification task than the same model trained on unfiltered web data, using 5-fold cross-validation. This extrinsic criterion aligns with \citet{xiao2023evaluating}'s concept of concurrent validity, substituting task performance for human expert ratings as the reference standard. Together, these three criteria provide a layered definition that moves from intrinsic linguistic properties to practical task performance, ensuring the corpus is not only clean by internal metrics but also useful for real NLP applications.


\subsection{Existing automated pipeline approaches}
The literature on automated web corpus construction reveals a clear division between two incompatible paradigms: generic multilingual pipelines that prioritize scale at the expense of linguistic accuracy, and language-specific manual pipelines that achieve quality but cannot be scaled. Neither adequately addresses the needs of low-resource African languages such as Sesotho, which require both automation and linguistically informed processing. Generic pipelines such as the Brno Corpus Processing Pipeline \citet{suchomel2020better} and the Colossal Clean Crawled Corpus \citep{dodge2021documenting} apply standardized, language-agnostic filters—n-gram models, blocklists, and statistical language identification—that treat all languages as statistically equivalent. \citet{suchomel2020better} acknowledges that yield rates vary significantly across languages but does not address how morphological complexity affects these rates, while \citet{dodge2021documenting} document that blocklist filtering disproportionately removes minority dialects and that langdetect performs poorly for non-dominant languages. The Labadian Crawler \citet{de2024data} follows a similar statistical approaches, with \citet{de2024data} reporting that their Tetun LID model exhibits biases toward frequent terms and achieves perfect tokenization only on texts conforming to a single standard orthography. None of these pipelines perform morphological analysis or adjust preprocessing based on writing system type. At the opposite extreme, language-specific manual pipelines such as the Swahili corpus \citet{masua2024heart} and the Tetun corpus \citet{de2024data} achieve high quality through intensive, custom effort—manually identifying sources, developing language-specific tokenizers from scratch, and training dedicated LID models—but these methods do not generalize. The CHB-EDC system \citet{Zhou2024} achieves high accuracy for structured medical form extraction but is not designed for general web corpus construction. \citet{nchabeleng2020evaluating} directly identify the specific failures of generic pipelines for Bantu languages: conjunctively versus disjunctively written languages require different preprocessing (morphological analysis versus part-of-speech tagging); stemming verbs and adjectives is crucial to avoid sparsity, but stemming nouns destroys semantic information encoded in noun class prefixes; threshold values in topic modeling must be adjusted for agglutination; and sentiment analysis should be performed before negation morphemes are lost. Their evaluation on Twitter data demonstrates that applying generic pipelines to Bantu languages produces systematically degraded results. For a language like Sesotho—a disjunctively written Bantu language with noun class systems, extensive concordial agreement, and thousands of possible verb forms per root—a generic pipeline would treat each inflected verb form as a separate token, fail to recognize negation morphemes as grammatical markers, and mishandle disjunctive writing. Conversely, replicating the Swahili or Tetun manual approach would require developing a custom tokenizer, part-of-speech tagger, and LID model from scratch. The literature has established robust crawling infrastructure \citep{suchomel2020better, de2024data} and separately documented the linguistic preprocessing requirements for Bantu languages \citet{nchabeleng2020evaluating}. However, no existing pipeline integrates domain-general crawling efficiency with language-specific morphological processing. The gap is not in automation technology or in linguistic knowledge but in their integration.

\subsection{Language identification as a pipeline component}
Language identification (LID) is a critical upstream component of web corpus construction, yet the literature indicates that LID for low-resource languages remains an unsolved problem. Existing tools and evaluation paradigms systematically fail to account for a fundamental asymmetry: false positives (misclassifying non-target text as a target language) are more damaging than false negatives because incorrectly included text cannot be reliably identified or removed after collection, whereas excluded text can potentially be recovered through additional crawling.

\citet{kreutzer-etal-2022-quality} demonstrate that standard evaluation metrics such as F1-score are misleading when applied to real-world web data. A model achieving 98\% F1 on clean benchmarks produced monolingual corpora with a median precision of only 5\% when evaluated on actual web crawls, due to extreme class imbalance. They recommend prioritising the false-positive rate (FPR) over F1 in real-world scenarios. \citet{kargaran2024glotcc} address this directly in GlotLID, introducing "und" labels for unseen scripts and "zxx" labels for six types of web noise (mis-rendered PDFs, Mojibake, binary files, etc.). Their design explicitly prioritises FPR reduction over raw F1: on 2,102 labels, GlotLID achieves 0.991 F1 but an FPR of just 0.000003. This low FPR ensures that text from other languages or noise is rarely misclassified as a target low-resource language.

The consequences of inadequate LID are severe for African languages. \citet{adebara2022afrolid} audit existing web-crawled datasets and find that at least 15 corpora were completely erroneous for African language content, 82 were mislabelled, and fastText embeddings for Yoruba contained 90\% words from other languages. They introduce AfroLID for 517 African languages, though as a Transformer model, it is less efficient for billion-scale processing.

\citet{fedorova2026openlid} provides recent production experience with OpenLID, evaluating multiple models on closely related languages. They find that OpenLID-v3 achieves better precision and FPR than its predecessor, while GlotLID achieves better recall; top-1 ensembling achieves the lowest FPR. Critically, they find that improving precision removes substantial data (e.g., 45\% of Venetian samples), demonstrating a real trade-off between precision and coverage. They conclude that reliable evaluation of LID in similar languages requires language-group-specific benchmarks rather than relying on high-level metrics.

For a corpus pipeline targeting a low-resource Bantu language like Sesotho, the implication is clear: false positives are more costly than false negatives. A document incorrectly classified as Sesotho irreversibly contaminates the final corpus; a document incorrectly rejected represents a recoverable loss. Therefore, LID selection should prioritise FPR minimisation over raw F1. GlotLID's explicit design philosophy of prioritising FPR \citet{kargaran2024glotcc} and its empirical performance with very low FPR make it suitable, subject to evaluation on language-group-specific benchmarks \citep{fedorova2026openlid}. The alternative—a higher-recall, higher-FPR model—risks producing a corpus whose Sesotho subset is contaminated with other languages and noise.

\section{Discussion}


\section{Summary}

\bibliographystyle{ruauthordate}
\bibliography{lit_review}  
\end{document}
